%------------------------------------------------------------------------------%
\begin{question}
  Diagonalize ortogonalmente a matriz
  \[
    A =
    \begin{pmatrix}
      5 & 2 \\
      2 & 2
    \end{pmatrix}
  \]
  Determine as matrizes $P$ e $D$ tais que
  \(
    A = P D P^t
  \)
\end{question}

%------------------------------------------------------------------------------%
\begin{resolution}{Autovalores}
\begin{columns}
\column{0.4\textwidth}
  \onslide<+->{Polinômio característico}
  \begin{align*}
    \onslide<+->{p(\lambda)}
    \onslide<+->{& = \det(A-\lambda I)}
    \\
    \onslide<+->{& = \begin{vmatrix}[c]
      5-\lambda & 2         \\
      2         & 2-\lambda
    \end{vmatrix}}
    \\
  \onslide<+->{  & = (5-\lambda)(2-\lambda) - 4 \\}
    \onslide<+->{& = \lambda^{2} - 7\lambda + 6}
  \end{align*}
\column{0.6\textwidth}
  \[
    \onslide<+->{\Delta}
    \onslide<+->{= b^2 - 4ac}
    \onslide<+->{= (-7)^2 - 4 \times 1 \times 6}
    \onslide<+->{= 25}
  \]

  \smallskip
  \[
    \onslide<+->{\sqrt{\Delta}}
    \onslide<+->{= \sqrt{25}}
    \onslide<+->{= 5}
  \]

  \smallskip
  \[
    \onslide<+->{\lambda_1}
    \onslide<+->{= \frac{-b-\sqrt{\Delta}}{2a}}
    \onslide<+->{= \frac{7-5}{2}}
    \onslide<+->{= 1}
  \]

  \smallskip
  \[
    \onslide<+->{\lambda_2}
    \onslide<+->{= \frac{-b+\sqrt{\Delta}}{2a}}
    \onslide<+->{= \frac{7+5}{2}}
    \onslide<+->{= 6}
  \]
\end{columns}
\end{resolution}

%------------------------------------------------------------------------------%
\begin{resolution}{Autovetores para $\lambda=1$}
\begin{columns}
\column{0.4\textwidth}
  \begin{align*}
    \onslide<+->{(A-I)\vec{x}&=\vec{0}}
    \\
    \onslide<+->{\begin{bmatrix}
      4 & 2 \\
      2 & 1
    \end{bmatrix}
    \vetor{x}{y} & = \vetor{0}{0}}
  \end{align*}
  \onslide<+->{\begin{align*}
    4x + 2y & = 0 \\
    2x +  y  & = 0
  \end{align*}}
\column{0.4\textwidth}
  \onslide<+->{Escolhendo $x=1$}
  \onslide<+->{\[
    y = -2
  \]}
  \onslide<+->{\[
    \vec{v}_1 = \vetor{1}{-2}
  \]}
  \onslide<+->{\[
    E(1) = \ger \left\{ \vetor{1}{-2} \right\}
  \]}
\end{columns}
\end{resolution}

%------------------------------------------------------------------------------%
\begin{resolution}{Normalização de $\vec{v}_1$}
  Norma de
  \(
    \vec{v}_1 = \vetor{1}{-2}
  \)
  \pause
  \[
    \norm{\vec{v}_{1}}
    \pause
    = \sqrt{1^{2}+(-2)^{2}}
    \pause
    = \sqrt{5}
  \]

  \pause
  Autovetor unitário
  \[
    \pause
    \vec{u}_{1}
    \pause
    = \frac{1}{\sqrt{5}}
      \vetor{1}{-2}
  \]
\end{resolution}

%------------------------------------------------------------------------------%
\begin{resolution}{Autovetores para $\lambda=6$}
\begin{columns}
\column{0.4\textwidth}
  \begin{align*}
    \onslide<+->{(A-6I)\vec{x}&=\vec{0}}
    \\
    \onslide<+->{\begin{bmatrix}
      -1 & 2  \\
      2  & -4
    \end{bmatrix}
    \vetor{x}{y} & = \vetor{0}{0}}
  \end{align*}
  \begin{align*}
    \onslide<+->{-x + 2y & = 0 \\}
    \onslide<+->{ x      & = 2y}
  \end{align*}
\column{0.4\textwidth}
  \onslide<+->{Escolhendo $y=1$}
  \onslide<+->{\[
    x = 2
  \]}
  \onslide<+->{\[
    \vec{v}_2 = \vetor{2}{1}
  \]}
  \onslide<+->{\[
    E(6) = \ger \left\{ \vetor{2}{1} \right\}
  \]}
\end{columns}
\end{resolution}

%------------------------------------------------------------------------------%
\begin{resolution}{Normalização de $\vec{v}_2$}
  Norma de
  \(
    \vec{v}_2 = \vetor{2}{1}
  \)

  \pause
  \[
    \norm{\vec{v}_{2}}
    \pause
    = \sqrt{2^{2}+1^{2}}
    \pause
    = \sqrt{5}
  \]

  \pause
  Autovetor unitário
  \[
    \pause
    \vec{u}_{2}
    \pause
    = \frac{1}{\sqrt{5}}
      \vetor{2}{1}
  \]
\end{resolution}

%------------------------------------------------------------------------------%
\begin{resolution}{Diagonalização}
  Organizamos os autovetores unitário nas colunas de $P$
  \[
    P =
    \frac{1}{\sqrt{5}}
    \begin{pmatrix}
      1  & 2 \\
      -2 & 1
    \end{pmatrix}
  \]

  \pause
  A matriz diagonal correspondente é
  \[
    D =
    \begin{pmatrix}
        1 & 0\\
        0 & 6
    \end{pmatrix}
  \]
\end{resolution}

%------------------------------------------------------------------------------%
